% !TeX program = xelatex
\documentclass[a4paper,10pt]{article}

% ============================================================================
% សន្លឹកកិច្ចការសិស្ស ស-ខ្មៅ (Student Worksheet - Grade 8 3.2.3 & 3.2.4 PCT2)
% បរិបទ៖ គ្មានស្លាយ LCD និងគ្មានម៉ាស៊ីនគិតលេខ — សិស្សបំពេញផ្ទាល់លើសន្លឹកកិច្ចការ ២ ទំព័រ A4
% ============================================================================
\usepackage[a4paper, left=0.85cm, right=0.85cm, top=0.70cm, bottom=0.95cm, footskip=0.48cm]{geometry}
\usepackage{amsmath, amssymb}
\usepackage{unicode-math}
\usepackage{fontspec}
\usepackage{array, colortbl}
\usepackage[table]{xcolor}
\usepackage{enumitem}
\usepackage{pifont}
\usepackage{fancyhdr}
\usepackage{lastpage}
\usepackage{tikz}
\usepackage[most]{tcolorbox}
\usetikzlibrary{calc, arrows.meta, decorations.pathreplacing, positioning, babel}

\makeatletter
\let\@expl@@@mark@update@singlecol@structures@@\relax
\makeatother

\providecommand{\NGSFindFonts}{%
  \IfFileExists{C:/Windows/Fonts/cambria.ttc}{%
    \def\NGSSysPath{C:/Windows/Fonts/}%
    \def\NGSUserPath{C:/Users/KIT/AppData/Local/Microsoft/Windows/Fonts/}%
  }{%
    \IfFileExists{./fonts/KhmerOSsiemreap.ttf}{\def\NGSSysPath{./fonts/}\def\NGSUserPath{./fonts/}}{%
    \IfFileExists{../fonts/KhmerOSsiemreap.ttf}{\def\NGSSysPath{../fonts/}\def\NGSUserPath{../fonts/}}{%
    \IfFileExists{../../fonts/KhmerOSsiemreap.ttf}{\def\NGSSysPath{../../fonts/}\def\NGSUserPath{../../fonts/}}{%
    \IfFileExists{../../../fonts/KhmerOSsiemreap.ttf}{\def\NGSSysPath{../../../fonts/}\def\NGSUserPath{../../../fonts/}}{%
    \IfFileExists{../../../../fonts/KhmerOSsiemreap.ttf}{\def\NGSSysPath{../../../../fonts/}\def\NGSUserPath{../../../../fonts/}}{%
    \IfFileExists{../../../../../fonts/KhmerOSsiemreap.ttf}{\def\NGSSysPath{../../../../../fonts/}\def\NGSUserPath{../../../../../fonts/}}{%
    \IfFileExists{../../../../../../fonts/KhmerOSsiemreap.ttf}{\def\NGSSysPath{../../../../../../fonts/}\def\NGSUserPath{../../../../../../fonts/}}{%
      \def\NGSSysPath{../../../../../../../fonts/}\def\NGSUserPath{../../../../../../../fonts/}%
    }}}}}}}%
  }%
}
\NGSFindFonts

\setmathfont[Path={\NGSSysPath}, FontIndex=1, Scale=1.02]{cambria.ttc}
\defaultfontfeatures{Ligatures=TeX}

\setmainfont[
  Path={\NGSUserPath}, Extension=.ttf, Script=Khmer,
  BoldFont={KhmerOSsiemreap}, BoldFeatures={FakeBold=1.8},
  ItalicFont={KhmerOSsiemreap}, ItalicFeatures={FakeSlant=0.22}
]{KhmerOSsiemreap}

\newfontfamily\muol[Path={\NGSUserPath}, Extension=.ttf, Script=Khmer]{KhmerOSmuollight}
\newfontfamily\TimesEn[Path={\NGSSysPath}, Extension=.ttf, BoldFont={timesbd}, ItalicFont={timesi}]{times}

\XeTeXlinebreaklocale "km"
\XeTeXlinebreakskip = 0pt plus 1pt minus 0.1pt

\definecolor{panelgray}{gray}{0.93}
\definecolor{headgray}{gray}{0.85}

\pagestyle{fancy}
\fancyhf{}
\renewcommand{\headrulewidth}{0pt}
\fancyfoot[L]{\fontsize{7.5pt}{9.5pt}\selectfont អនុវិទ្យាល័យព្រះក្រសាំង \en{|} គណិតវិទ្យាថ្នាក់ទី៨ \en{|} មេរៀនទី៣៖ ២.៣ ការបញ្ចុះតម្លៃ និង ២.៤ បម្រែបម្រួលភាគរយ (គ្រូ៖ ហុក ប៊ុនថែន)}
\fancyfoot[R]{\fontsize{7.8pt}{9.5pt}\selectfont ទំព័រ \en{\thepage/\pageref{LastPage}}}

\setlength{\parindent}{0pt}
\setlength{\parskip}{0pt}
\setlength{\tabcolsep}{3.8pt}

\newcommand{\en}[1]{{\TimesEn #1}}
\newcommand{\enbf}[1]{{\TimesEn\bfseries #1}}
\newcommand{\unit}[1]{\,\text{\en{#1}}}
\newcommand{\cnum}[1]{\mbox{\fontsize{10pt}{12pt}\selectfont\ding{\numexpr181+#1\relax}}}
\newcommand{\chkbox}{\tikz[baseline=-0.12ex]\draw[black, line width=0.8pt] (0,0) rectangle (0.24cm,0.24cm);}
\newcommand{\blank}[1][1.5cm]{\makebox[#1]{\dotfill}}
\newcommand{\ansbox}[1][0.85cm]{\,\tikz[baseline=-0.5ex]\draw[black, line width=0.65pt, fill=white] (0,0) rectangle (#1, 0.42cm);\,}
\newcommand{\secbar}[2]{%
  \begin{tcolorbox}[enhanced, sharp corners, colback=headgray, colframe=black, boxrule=0.9pt, boxsep=1pt, top=2pt, bottom=2pt, left=5pt, right=5pt, before skip=4pt, after skip=3pt]
    {\muol\fontsize{9.2pt}{13pt}\selectfont #1}\hfill{\fontsize{8pt}{11pt}\selectfont\textbf{#2}}
  \end{tcolorbox}%
}

\newcolumntype{L}[1]{>{\raggedright\arraybackslash}p{#1}}
\newcolumntype{C}[1]{>{\centering\arraybackslash}p{#1}}

\begin{document}
\fontsize{9.6pt}{14.4pt}\selectfont

% ============================================================================
% ទំព័រទី ១៖ E1 + E2 (POE 1 & POE 2) + E3 សំយោគរូបមន្ត និងធនាគារហេតុផល
% ============================================================================
\begin{tcolorbox}[enhanced, sharp corners, colback=white, colframe=black, boxrule=1.0pt, boxsep=2pt, top=3pt, bottom=3pt, left=6pt, right=6pt]
  \begin{center}
    {\muol\fontsize{10.0pt}{15pt}\selectfont អនុវិទ្យាល័យព្រះក្រសាំង \enbf{--} សន្លឹកកិច្ចការសិស្ស \enbf{(}មេរៀនទី៣៖ ២\enbf{.}៣\enbf{--}២\enbf{.}៤ ការបញ្ចុះតម្លៃ និងបម្រែបម្រួលភាគរយ\enbf{)}}\\[3pt]
    {\fontsize{8.8pt}{12pt}\selectfont
    ឈ្មោះសិស្ស / ដៃគូ៖ \blank[5.2cm] \hspace{0.25cm}
    ថ្នាក់ទី៨៖ \blank[1.1cm] \hspace{0.25cm}
    ក្រុមទី៖ \blank[1.0cm] \hspace{0.25cm}
    កាលបរិច្ឆេទ៖ ៣០/១១/២០២៦}
  \end{center}
\end{tcolorbox}

% --- E1 + E2 ---
\secbar{ផ្នែកទី១៖ បេសកកម្មទស្សន៍ទាយ \enbf{(E1)} និងការរុករកជាដៃគូ \enbf{(E2 -- POE 1 \& POE 2)}}{រយៈពេល ១៤ នាទី}

\begin{tcolorbox}[enhanced, sharp corners, colback=panelgray, colframe=black, boxrule=0.75pt, boxsep=1.5pt, top=2.5pt, bottom=2.5pt, left=5pt, right=5pt]
  \textbf{បេសកកម្ម \en{E1} (គូសធីក \chkbox\ មុនពិសោធន៍)៖}
  (១) កាបូបដាក់លក់ថ្លៃ $32\,000\text{៛}$ លក់ជាក់ស្ដែង $25\,600\text{៛}$ — តើបញ្ចុះតម្លៃ៖ \chkbox\ $20\%$ \quad \chkbox\ $25\%$\\
  (២) ចន្ធីទទួលបានប្រាក់ខែថ្មី $518\,400\text{៛}$ ក្រោយតម្លើង $8\%$ — តើប្រាក់ខែដើមស្មើ៖ \chkbox\ $476\,928\text{៛}$ \quad \chkbox\ $480\,000\text{៛}$
\end{tcolorbox}

\vspace{3pt}
\textbf{ក. សកម្មភាពទី ១ (\en{POE 1})៖ យុថ្កាថ្លៃដាក់លក់ $100\%$ លើការបញ្ចុះតម្លៃ (សៀវភៅពុម្ព ទំ.៣៣)}

{\renewcommand{\arraystretch}{1.46}
\begin{tabular}{|L{4.50cm}|C{2.90cm}|C{2.90cm}|C{3.50cm}|C{3.80cm}|}
  \hline
  \rowcolor{headgray}
  \textbf{មុខទំនិញ} & \textbf{ថ្លៃដាក់លក់ ($100\%$)} & \textbf{ប្រាក់បញ្ចុះតម្លៃ} & \textbf{ថ្លៃលក់ជាក់ស្ដែង} & \textbf{ភាគរយបញ្ចុះតម្លៃ} \\
  \hline
  កាបូបយួរដៃ (ទំ.៣៣) & $32\,000\text{៛}$ & $32\,000-25\,600=\blank[0.9cm]$ & $25\,600\text{៛}$ & $\dfrac{6400}{\blank[1.1cm]}\times 100\% = \blank[0.7cm]\%$ \\
  \hline
  អាវធំបញ្ចុះ $30\%$ (ទំ.៣៣) & $36\,000\text{៛}$ & $36\,000\times\frac{30}{100}=\blank[1.0cm]$ & $36\,000-\blank[1.0cm]=\blank[1.0cm]$ & បញ្ចុះ $30\%$ (សល់ $70\%$) \\
  \hline
\end{tabular}}

\vspace{3pt}
$\bullet$ \textbf{ពន្យល់ (\en{E})៖} ភាគរយនៃការបញ្ចុះតម្លៃ ត្រូវចែកនឹង \blank[2.3cm] ($100\%$) ជានិច្ច!

\vspace{4pt}
\textbf{ខ. សកម្មភាពទី ២ (\en{POE 2})៖ បន្ទះមាត្រដ្ឋានភ្លោះ — ទៅមុខ (ប៊ុនថា) ទល់នឹង ត្រឡប់ក្រោយ (ចន្ធី) (ទំ.៣៤–៣៥)}

{\renewcommand{\arraystretch}{1.46}
\begin{tabular}{|L{4.50cm}|C{3.10cm}|C{3.30cm}|L{6.70cm}|}
  \hline
  \rowcolor{headgray}
  \textbf{ករណីប្រាក់ខែ} & \textbf{ប្រាក់ខែដើម ($100\%$)} & \textbf{ប្រាក់ខែថ្មី ($108\%$)} & \textbf{សមាមាត្រគណនាដោយដៃ} \\
  \hline
  ១) ប៊ុនថា (ឡើង $+8\%$) & $540\,000\text{៛}$ (ស្គាល់) & $\blank[1.8cm]\text{៛}$ & $\text{ថ្មី} = 540\,000 \times \dfrac{108}{100} = \blank[1.8cm]\text{៛}$ \\
  \hline
  ២) ចន្ធី (ឡើង $+8\%$) & $x = \blank[1.5cm]\text{៛}$ & $518\,400\text{៛}$ (ស្គាល់) & $\dfrac{x}{100} = \dfrac{518\,400}{\blank[0.9cm]} \Rightarrow x = \blank[1.8cm]\text{៛}$ \\
  \hline
\end{tabular}}

\vspace{3pt}
$\bullet$ \textbf{ពន្យល់ (\en{E})៖} ហេតុអ្វីយក $518\,400\text{៛}$ ដក $8\%$ បាន $476\,928\text{៛}$ ខុស? ព្រោះ $8\%$ នោះគិតលើ \blank[2.3cm] មិនមែនលើ $518\,400\text{៛}$ ទេ!

% --- E3 ---
\secbar{ផ្នែកទី២៖ សំយោគរូបមន្តគោល \enbf{(E3)} និងធនាគារហេតុផល \enbf{B\&W} \cnum{1}\enbf{--}\cnum{5}}{សៀវភៅពុម្ព ទំ.៣៣–៣៥}

\begin{tabular}{|L{9.20cm}|L{9.20cm}|}
  \hline
  \textbf{១. រូបមន្តបញ្ចុះតម្លៃ និងបម្រែបម្រួលភាគរយ៖}\par\vspace{2pt}
  $\bullet$ \textbf{ភាគរយបញ្ចុះតម្លៃ៖} $\dfrac{\text{ប្រាក់បញ្ចុះតម្លៃ}}{\blank[2.2cm]} \times 100\%$\par\vspace{2.5pt}
  $\bullet$ \textbf{ថ្លៃលក់ក្រោយបញ្ចុះ $d\%$៖} $\text{ថ្លៃដាក់លក់} \times \dfrac{100 - \blank[0.7cm]}{100}$\par\vspace{2.5pt}
  $\bullet$ \textbf{បម្រែបម្រួលទៅមុខ ($\pm p\%$)៖} $\text{តម្លៃថ្មី} = \text{ដើម} \times \dfrac{100 \pm p}{100}$\par\vspace{2.5pt}
  $\bullet$ \textbf{រកតម្លៃដើមត្រឡប់ក្រោយ៖} $\text{ដើម} = \text{តម្លៃថ្មី} \times \dfrac{100}{\blank[1.5cm]}$
  &
  \textbf{២. ធនាគារហេតុផល \en{B\&W} (\cnum{1}--\cnum{5})៖}\par\vspace{2pt}
  \cnum{1} យុថ្កាបញ្ចុះតម្លៃ៖ ថ្លៃដាក់លក់ស្មើ $100\%$ ជាភាគបែងជានិច្ច ។\par\vspace{2pt}
  \cnum{2} ញែកយុថ្កា៖ បញ្ចុះតម្លៃធៀបនឹងថ្លៃដាក់លក់ ; ចំណេញធៀបនឹងថ្លៃដើម ។\par\vspace{2pt}
  \cnum{3} មេគុណកើន-ថយ៖ កើន $p\% \to (100+p)\%$ ; ថយ $d\% \to (100-d)\%$ ។\par\vspace{2pt}
  \cnum{4} រកដើមក្រោយកើន $p\%$៖ ចែកតម្លៃថ្មីនឹង $(100+p)\%$ រួចគុណ $100$ ។\par\vspace{2pt}
  \cnum{5} រកថ្លៃដាក់លក់ក្រោយបញ្ចុះ $d\%$៖ $\text{ថ្លៃដាក់លក់} = \frac{100}{100-d}\times\text{ថ្លៃលក់}$ ។ \\
  \hline
\end{tabular}

\vspace{3pt}
\begin{tcolorbox}[enhanced, sharp corners, colback=panelgray, colframe=black, boxrule=0.65pt, boxsep=1.5pt, top=2pt, bottom=2pt, left=5pt, right=5pt]
  \textbf{ប្រយ័ត្នអន្ទាក់៖} ហាមយកតម្លៃថ្មីទៅដក $p\%$ ឬបូក $d\%$ ដើម្បីរកតម្លៃដើម! ត្រូវសរសេរសមាមាត្រធៀបនឹង $(100\pm p)\%$ ជានិច្ច!
\end{tcolorbox}

\vspace{3pt}
\textbf{៣. លំហាត់ត្រួតពិនិត្យរហ័ស៖ គណនាការបញ្ចុះតម្លៃ និងបម្រែបម្រួលភាគរយ}

\vspace{2pt}
{\renewcommand{\arraystretch}{1.60}
\begin{tabular}{|C{4.40cm}|C{4.40cm}|C{4.40cm}|C{4.40cm}|}
  \hline
  \rowcolor{panelgray}
  ដាក់លក់ $50\$$ បញ្ចុះ $10\% \to \blank[0.8cm]$ & ដាក់លក់ $80\$$ លក់ $60\$$៖ $-\blank[0.7cm]\%$ & ដើម $200\$$ កើន $15\% \to \blank[0.8cm]$ & ដើម $300\$$ ថយ $20\% \to \blank[0.8cm]$ \\
  \hline
  កើន $10\%$ បាន $220\$$៖ ដើម $\blank[0.8cm]$ & បញ្ចុះ $20\%$ សល់ $40\$$៖ ដើម $\blank[0.7cm]$ & ដាក់លក់ $40\,000\text{៛}$ បញ្ចុះ $25\% \to \blank[0.7cm]$ & កើន $25\%$ បាន $100\$$៖ ដើម $\blank[0.7cm]$ \\
  \hline
  បញ្ចុះ $10\%$ សល់ $90\$$៖ ដើម $\blank[0.7cm]$ & ដាក់លក់ $120\$$ លក់ $90\$$៖ $-\blank[0.6cm]\%$ & ដើម $60\,000\text{៛}$ ឡើង $5\% \to \blank[0.8cm]$ & បញ្ចុះ $50\%$ សល់ $35\$$៖ ដើម $\blank[0.7cm]$ \\
  \hline
  ឡើង $20\%$ បាន $600\$$៖ ដើម $\blank[0.7cm]$ & ដាក់លក់ $25\,000\text{៛}$ បញ្ចុះ $20\% \to \blank[0.6cm]$ & ដើម $400\$$ ថយ $12\% \to \blank[0.8cm]$ & បញ្ចុះ $25\%$ សល់ $75\$$៖ ដើម $\blank[0.7cm]$ \\
  \hline
\end{tabular}}

\newpage

% ============================================================================
% ទំព័រទី ២៖ E4 លំហាត់គំរូ + E5 ប្រតិបត្តិ + វិនិច្ឆ័យ A-D + Exit Check + កិច្ចការផ្ទះ
% ============================================================================
\secbar{ផ្នែកទី៣៖ ពង្រីកចំណេះដឹងលើចំណោទអាជីវកម្មចម្រុះ \enbf{(E4 --} ទំ\enbf{.}៣៤\enbf{)}}{រយៈពេល ១០ នាទី}

\textbf{លំហាត់គំរូ (សៀវភៅពុម្ព ទំ.៣៤ — មីងសារីលក់ឌីស \en{CD})៖} មីងសារីទិញឌីស \en{CD} ចំនួន $\mathbf{300}$ បន្ទះ ដោយក្នុង $1$ បន្ទះថ្លៃដើម $\mathbf{3300\text{៛}}$ ។ គាត់ដាក់លក់ក្នុង $1$ បន្ទះថ្លៃ $\mathbf{3980\text{៛}}$ ហើយលក់ដាច់បាន $\mathbf{200}$ បន្ទះ ។ ឌីសដែលនៅសល់ គាត់បញ្ចុះតម្លៃ $\mathbf{15\%}$ ហើយលក់ដាច់អស់ ។ រកភាគរយប្រាក់ចំណេញសរុបរបស់មីងសារី៖

{\renewcommand{\arraystretch}{1.48}
\begin{tabular}{|L{4.60cm}|L{9.00cm}|C{4.20cm}|}
  \hline
  \rowcolor{headgray}
  \textbf{ជំហានគណនា} & \textbf{ការដាក់រូបមន្ត និងគណនាដោយដៃ} & \textbf{លទ្ធផល និងហេតុផល} \\
  \hline
  \textbf{១. ប្រាក់ថ្លៃដើមសរុប ($300$ បន្ទះ)} & $\text{ថ្លៃដើមសរុប} = 300 \times 3300 = \blank[1.8cm]\text{៛}$ & យុថ្កាថ្លៃដើម $100\%$ (\cnum{2}) \\
  \hline
  \textbf{២. ប្រាក់លក់ $200$ បន្ទះដំបូង} & $200 \times 3980 = \blank[1.8cm]\text{៛}$ & លក់តាមថ្លៃដាក់លក់ \\
  \hline
  \textbf{៣. ប្រាក់លក់ $100$ បន្ទះបញ្ចុះ $15\%$} & $1$ បន្ទះ $= 3980\times\frac{85}{100} = \blank[1.2cm]\text{៛} \implies 100$ បន្ទះ $= \blank[1.5cm]\text{៛}$ & បញ្ចុះលើ $3980\text{៛}$ (\cnum{1}) \\
  \hline
  \textbf{៤. ប្រាក់លក់សរុប និងចំណេញ $\%$} & លក់សរុប $= \blank[1.6cm]\text{៛} \implies \text{ចំណេញ}\% = \dfrac{144\,300}{\blank[1.4cm]}\times 100\% \approx \blank[1.0cm]\%$ & ចំណេញ $\approx \blank[1.0cm]\%$ (\cnum{2}) \\
  \hline
\end{tabular}}

\vspace{4pt}
$\bullet$ \textbf{សំណួរឆ្លុះបញ្ចាំង \en{(EIA Looking-Back)}៖} បើទំនិញមួយតម្លើងថ្លៃ $20\%$ រួចបញ្ចុះ $20\%$ វិញ តើស្មើថ្លៃដើម $100\%$ វិញឬទេ? \blank[2.8cm]

% --- E5 ---
\secbar{ផ្នែកទី៤៖ វាយតម្លៃលទ្ធផលសិក្សា \enbf{(E5 --} ប្រតិបត្តិ ទំ\enbf{.}៣៥ \enbf{+} វិនិច្ឆ័យ \enbf{A--D +} សំបុត្រចេញ\enbf{)}}{រយៈពេល ១០ នាទី}

\textbf{ក. លំហាត់ប្រតិបត្តិ (សៀវភៅពុម្ព ទំ.៣៤–៣៥ — ធ្វើជាដៃគូ រួចប្ដូរសន្លឹកផ្ទៀង \en{Pause-Review})៖}

{\renewcommand{\arraystretch}{1.48}
\begin{tabular}{|L{8.90cm}|L{8.90cm}|}
  \hline
  \textbf{១) ប្រតិបត្តិ (ទំ.៣៥ — រកថ្លៃដាក់លក់កង់មុនបញ្ចុះ $15\%$)៖}\par\vspace{2pt}
  កង់មួយបញ្ចុះតម្លៃ $\mathbf{15\%}$ (នៅសល់ $\mathbf{85\%}$) លក់ក្នុងថ្លៃកែតម្រូវគត់ $\mathbf{816\,000\text{៛}}$ \textit{(លេខសៀវភៅពុម្ព $846\,000\text{៛}$)} ។ រកថ្លៃដើមមុនបញ្ចុះ៖\par\vspace{2pt}
  $\bullet$ សមាមាត្រ៖ $\dfrac{x}{100} = \dfrac{816\,000}{\blank[1.0cm]}$ (\cnum{5})\par\vspace{2pt}
  $\bullet$ $x = \dfrac{100}{85}\times 816\,000 = \dfrac{20}{17}\times 816\,000 = \blank[1.8cm]\text{៛}$\par\vspace{2pt}
  \textit{(បើគណនាតាម $846\,000\text{៛}$៖ $x = \frac{100}{85}\times 846\,000 \approx 995\,294\text{៛}$)}
  &
  \textbf{២) ពង្រឹង (ទំ.៣៤ — បម្រែបម្រួលប្រាក់ខែកម្មករ)៖}\par\vspace{2pt}
  កម្មកររោងចក្រម្នាក់មានប្រាក់ខែដើម $\mathbf{320\,000\text{៛}}$ ។ គណនាប្រាក់ខែថ្មីរបស់គាត់ប្រសិនបើ៖\par\vspace{2pt}
  $\bullet$ \textbf{ក. តម្លើង $15\%$៖} $320\,000 \times \dfrac{\blank[0.9cm]}{100} = \blank[1.7cm]\text{៛}$ (\cnum{3})\par\vspace{2pt}
  $\bullet$ \textbf{ខ. បន្ថយ $5\%$៖} $320\,000 \times \dfrac{\blank[0.9cm]}{100} = \blank[1.7cm]\text{៛}$ \\
  \hline
\end{tabular}}

\vspace{5pt}
\textbf{ខ. គូសវង់ជុំវិញជម្រើសត្រឹមត្រូវតែ ១ គត់ ក្នុងចំណោម \en{A, B, C, D} រួចសរសេរលេខហេតុផល \cnum{1}--\cnum{5} នៃជម្រើសខុស៖}

{\renewcommand{\arraystretch}{1.36}
\begin{tabular}{|C{4.40cm}|C{4.40cm}|C{4.40cm}|C{4.40cm}|}
  \hline
  \cellcolor{panelgray}\textbf{ជម្រើស \en{A}} & \cellcolor{panelgray}\textbf{ជម្រើស \en{B}} & \cellcolor{panelgray}\textbf{ជម្រើស \en{C}} & \cellcolor{panelgray}\textbf{ជម្រើស \en{D}} \\
  \hline
  ចន្ធីឡើង $8\%$ បាន $518\,400\text{៛}$ $\Rightarrow$ ដើម $480\,000\text{៛}$ & ចន្ធី៖ $518\,400 - 8\%$ $= 476\,928\text{៛}$ & កាបូប $32\,000\to 25\,600$ បញ្ចុះតម្លៃ $25\%$ & កង់បញ្ចុះ $15\%$ សល់ $816\,000 \Rightarrow 938\,400\text{៛}$ \\[3pt]
  \chkbox\ ត្រូវ \quad \chkbox\ ខុស (\blank[0.7cm]) & \chkbox\ ត្រូវ \quad \chkbox\ ខុស (\blank[0.7cm]) & \chkbox\ ត្រូវ \quad \chkbox\ ខុស (\blank[0.7cm]) & \chkbox\ ត្រូវ \quad \chkbox\ ខុស (\blank[0.7cm]) \\
  \hline
\end{tabular}}

\vspace{5pt}
\begin{tcolorbox}[enhanced, sharp corners, colback=white, colframe=black, boxrule=0.85pt, boxsep=2.5pt, top=3.5pt, bottom=3.5pt, left=5pt, right=5pt]
  {\muol\fontsize{8.8pt}{12.5pt}\selectfont សំបុត្រចេញរហ័ស \enbf{(Exit Check --} ៦០ វិនាទី\enbf{)}៖}\\[2pt]
  (១) ស្បែកជើងដាក់លក់ $50\,000\text{៛}$ បញ្ចុះតម្លៃ $12\%$ លក់ថ្លៃ៖ \blank[1.8cm]៛ \quad\en{|}\quad
  (២) អាវបញ្ចុះ $20\%$ លក់សល់ $64\,000\text{៛}$ តើថ្លៃដាក់លក់ដើមប៉ុន្មាន? \blank[1.8cm]៛
\end{tcolorbox}

\vspace{3pt}
\begin{tcolorbox}[enhanced, sharp corners, colback=panelgray, colframe=black, boxrule=0.75pt, boxsep=2pt, top=3pt, bottom=3pt, left=5pt, right=5pt]
  {\muol\fontsize{8.8pt}{12.5pt}\selectfont កិច្ចការផ្ទះ \enbf{(}សៀវភៅគោល ទំព័រ ៤១\enbf{--}៤២ លំហាត់ \enbf{\#14} និង \enbf{\#15)}៖}\\[2pt]
  $\bullet$ \textbf{លំហាត់ \#14 (ទំ.៤១)៖} គណនាភាគរយបញ្ចុះតម្លៃ និងប្រាក់ចំណេញ-ខាតតាមយុថ្កាគោលនីមួយៗ ។\\
  $\bullet$ \textbf{លំហាត់ \#15 (ទំ.៤២)៖} គណនាបម្រែបម្រួលភាគរយ និងអានមុន \textbf{២.៥. ការប្រាក់សាមញ្ញ (ទំ.៣៥–៣៧)} ។
\end{tcolorbox}

\end{document}
